\subsection{Representation by ideal lattices}\label{app:ideal-representation}

The derived union gives a representation of every algebra in the varieties
generated by $\mathbf C$, $\mathbf A$, and $\mathbf A^{\rm sym}$, not just of
their finite generators. The carrier is recovered from an algebraic complete
lattice; owned choices become maps into lattices of biideals. We give the
order-theoretic construction directly. Its semilattice part is the ideal
form of semilattice duality; see \cite[Section~2]{RomanowskaSmith2017} for
the corresponding character-space treatment.

A join-semilattice below is nonempty and has no distinguished bottom.
Its ideals may be empty. Write $\mathcal I(B)$ for their inclusion-ordered
complete lattice. Joins are ideal generation of unions, and directed
nonempty joins are unions. For an algebraic complete lattice $L$, write
$K^+(L)$ for its nonzero compact elements. Here compact means that
$c\leq\bigvee D$ for a nonempty directed set $D$ implies $c\leq d$ for
some $d\in D$; in particular, bottom is compact and is explicitly excluded
from $K^+(L)$.

Call a map between complete lattices an \emph{ideal map} if it preserves
arbitrary meets, bottom, and nonempty directed joins. Preservation of the
empty meet includes preservation of top. Such a map $g:L\to M$ has a
left adjoint $g_!:M\to L$.

\begin{theorem}[Ideal representation]\label{olv:thm:ideal-representation}
Let $K\geq3$ and let $B$ be any algebra in the variety generated by one of
$\mathbf C$, $\mathbf A$, or $\mathbf A^{\rm sym}$. Equip $B$ with its
term-defined union $\vee$. Then $L=\mathcal I(B)$ is a nontrivial algebraic
complete lattice, and
\[
 B\longrightarrow K^+(L),\qquad b\longmapsto\mathord\downarrow b
\]
is a join-semilattice isomorphism onto all nonzero compact elements.
Each owned operation is recovered from the ideal map
\[
 \delta_r:L\longrightarrow\operatorname{Bi}(B,B),\qquad
 \delta_r(J)=\{(a,b):a\ow r b\in J\},
\]
where a biideal is a downward closed subset of $B^2$ closed under joins
separately in each coordinate. Specifically, writing
$P(a,b)=\mathord\downarrow a\times\mathord\downarrow b$,
\[
 (\delta_r)_!P(a,b)=\mathord\downarrow(a\ow r b).
\]
Winner constants are represented by their principal ideals. Each role
permutation is represented by inverse image on ideals, and its left
adjoint recovers its action on the nonzero compact elements. Thus the
representation is an isomorphism in the original signature in all three
cases.
\end{theorem}
\begin{proof}
Every nonempty ideal is the directed union of its principal subideals.
Each principal ideal is compact: membership of its generator in a directed
union already holds in one member. Conversely, compactness applied to
the union of principal subideals makes a nonempty compact ideal principal.
This proves algebraicity and the displayed carrier isomorphism.

Associativity, commutativity, and idempotence of the derived union, and
separate preservation of that union by each owned choice, are identities
of the generating winner algebra. They therefore hold in $B$. The
preimage defining $\delta_r(J)$ is a biideal: downward closure follows
from monotonicity, and the two join closures from separate preservation.
As an inverse-image map, $\delta_r$ preserves meets, bottom, and directed
unions. Finally,
\[
 P(a,b)\subseteq\delta_r(J)
 \quad\Longleftrightarrow\quad a\ow r b\in J
 \quad\Longleftrightarrow\quad
 \mathord\downarrow(a\ow r b)\subseteq J.
\]
The adjunction gives the recovery formula. Constants and permutations
are recovered by the same principal-ideal map; inverse image reverses
composition of the permutations.
\end{proof}

\paragraph{Recovery in the other direction.}
The construction is reversible and includes homomorphisms. If $L$ is a
nontrivial algebraic complete lattice, then
\[
 L\cong\mathcal I(K^+(L)),\qquad
 x\longmapsto\{c\in K^+(L):c\leq x\}.
\]
Indeed the supremum of this ideal is $x$ by algebraicity. Conversely, if
a positive compact lies below the supremum of a nonempty ideal, compactness
puts it below one member and hence in that ideal. The empty ideal gives
bottom.

An ideal map's left adjoint preserves joins and positive compactness.
For the latter, if $c$ is compact and $g_!c\leq\bigvee D$, the adjunction
and preservation of directed joins give
$c\leq g(\bigvee D)=\bigvee g[D]$; compactness supplies a witness in $D$.
Also $g_!c=0$ would imply $c\leq g(0)=0$.
Consequently inverse image gives an exact correspondence
\[
 \operatorname{Hom}_{\vee}(B,C)
 \cong\operatorname{IdealMap}(\mathcal I(C),\mathcal I(B)).
\]
For its converse, recover $h(b)$ from
$g_!(\mathord\downarrow b)=\mathord\downarrow h(b)$.
Then $b\in g(J)$ if and only if $h(b)\in J$, proving that $g=h^{-1}$.

For the operations, let
$\delta:L\to\operatorname{Bi}(K^+(L),K^+(L))$ be an ideal map. Define
\[
 m(c,d)=\delta_!P(c,d).
\]
Principal rectangles are positive compact biideals, so the result is in
$K^+(L)$. Joins of rectangles in the biideal lattice satisfy
$P(c\vee c',d)=P(c,d)\vee P(c',d)$, and likewise on the right.
Thus $m$ preserves joins separately. Its exact recovery identity is
\[
 (c,d)\in\delta(x)\quad\Longleftrightarrow\quad m(c,d)\leq x.
\]
These constructions are inverse on both carriers and operations.

The dual objects are therefore nontrivial algebraic complete lattices
with one such co-operation for each binary symbol, a positive compact
point for each constant, and a complete-lattice automorphism for each
distinguished semilattice automorphism. Extra equations, including any
group action and its compatibility with the other symbols, are imposed
through the recovery formula.

More explicitly, an ideal map $g:L\to M$ corresponds to the reverse
join homomorphism $h=K^+(g_!):K^+(M)\to K^+(L)$. Put
$\operatorname{Bi}(g)(R)=(h\times h)^{-1}[R]$. The homomorphism preserves
owned choices exactly when
\[
 \varepsilon_r\circ g=\operatorname{Bi}(g)\circ\delta_r.
\]
This follows by testing membership at $(c,d)$ and applying the recovery
identity and the adjunction. For constant points $a_i\in K^+(L)$ and
$b_i\in K^+(M)$, the condition is $h(b_i)=a_i$, equivalently
$b_i\leq g(x)\Longleftrightarrow a_i\leq x$ for every $x\in L$.
If $\alpha_p^L$ and $\alpha_p^M$ denote the distinguished inverse-image
automorphisms, the permutation condition is
$g\circ\alpha_p^L=\alpha_p^M\circ g$.
These are precisely the dual morphism conditions. Identity and composition
follow from inverse images. This is a contravariant equivalence for
join-semilattices with separately join-preserving operations, constants,
and distinguished semilattice automorphisms. Inverse image changes a
left permutation action into a right action.
It restricts to each winner variety by translating its identities through
the recovery formula. In making this restriction, union is a definitional
expansion: it must equal the fixed derived union term, rather than an
independent semilattice operation. The resulting dual description uses
the whole identity theory; it does not supply short intrinsic axioms for
the restricted class of ideal lattices.
